\lecture{3}{Mon 27 Jan 2025 13:25}{Rings with Unity}

Since every ring can be ``given'' unity, we generally assume rings have unity, because it makes life much easier.

\begin{definition}[Unit]\label{dfn:1.4}
	Let $R$ be a ring with unity, and let $a\in R$ such that there exists $b\in R$ such that $ab=1$. Then $a,b$ are units, and we say $b$ is the multiplicative inverse of $a$.
\end{definition}

\begin{proposition}\label{prp:1.2}
	Let $R$ be a ring with unity. Then
	\begin{itemize}
		\item the unity is unique;
		\item If an element $a\in R$ has a multiplicative inverse, then this inverse is unique.
	\end{itemize}
\end{proposition}
\begin{proof}
	Let $1_R,1_S$ be unities of $R$. Then $1_R1_S=1_R=1_S$ since they are both unities. Let $a\in R$ be a unit, and let $b,c$ be inverses. Then $1=ab=ac$, and it follows that
	\[
		bab=bac \implies 1b=1c \implies b=c
	.\]
\end{proof}
\begin{proposition}\label{prp:1.3}
	Let $R$ be a ring with unity. Then the set $U(R)$ of units of the ring $R$ is a group with respect to ring multiplication.
\end{proposition}
\begin{proof}
	Note that $1\in U(R)$, so clearly $U(R)$ has an identity. We also have that all elements have inverses by definition, and associativity follows from the ring axioms. Let $a,b\in U(R)$. Since $(ab)^{-1} =b^{-1} a^{-1} $, and $a^{-1} ,b^{-1} \in U(R)$, we have that $ab\in U(R)$, and $U(R)$ is a group.
\end{proof}
\begin{theorem}\label{thm:1.1}
	\[
		U(\mathbb{Z} _n)=\left\{ r\in\mathbb{Z} _n\mid \operatorname{gcd} (r,n)=1 \right\}
	.\]
\end{theorem}
\begin{proof}
	If the gcd is not 1, then there exists some multiple $mr=0$ of $r$ with $m<n$. If $r$ is a unit, then $mrr^{-1} =0$, so $m=0$, so $\operatorname{gcd} (r,n)=1$, which is a contradiction.
\end{proof}
Cancellation:
\begin{proposition}\label{prp:1.8}
	Let $R$ be a unital ring and let $a\in R$ be a unit. Then $ab=ac$ implies $b=c$.
\end{proposition}
\begin{proof}
	Let $ab=ac$. Then $a^{-1} ab=a^{-1} ac$, and thus $b=c$.
\end{proof}

An example of a finite field is $\mathbb{Z} _p$ for $p$ prime.
\begin{definition}[Subring]\label{dfn:1.4}
	Let $R$ be a ring and let $S\subseteq R$ be a subset. Then $S$ is a subring if it is a ring given the same operations as $R$; $S$ is a subring if the inclusion map $S\hookrightarrow R$ is a ring homomorphism.
\end{definition}
\begin{proposition}\label{prp:1.5}
	Let $S\subseteq R$, such that for all $a,b\in S$, $a-b\in S$ and $ab\in S$. Then $S$ is a subring of $R$.
\end{proposition}
\begin{proof}
	Trivial.
\end{proof}

